\documentclass[11pt]{article}

% Required packages:
\def\AssignmentMode{HW}     % HW | Lab | LabVIEW
\def\HWNum{Code Commenting}
\def\AssignmentText{All homework code is expected to be \textcolor{red}{commented} and have your \textcolor{red}{initialed} on each of the lines. Points WILL be taken off if this is not done.}

% Course metadata
% ============================
% Pre-load options
% ----------------------------
% xcolor must receive [table] BEFORE any package (tcolorbox, tikz, ...)
% pulls it in without options. Otherwise, "\usepackage[table]{xcolor}"
% later is a silent no-op under pdflatex but triggers an option-clash
% cascade under tex4ht ("Paragraph ended before \@reset@ptions").
% ============================
\PassOptionsToPackage{table}{xcolor}

% ============================
% Engine / output mode detection
% ----------------------------
% \ifPDFTeX  - true under pdflatex (also under tex4ht's htlatex driver)
% \ifLuaTeX  - true under lualatex
% \ifXeTeX   - true under xelatex
% \ifSEhtml  - true when being processed by tex4ht / make4ht
%             (named "SEhtml" to avoid clashing with package internals
%              in the @-letter namespace; the @ form was being reset
%              somewhere downstream, silently flipping the flag back.)
% ============================
\usepackage{iftex}

\newif\ifSEhtml
\SEhtmlfalse
\makeatletter
\@ifundefined{HCode}{}{\SEhtmltrue}
\makeatother

% ============================
% Geometry
% ============================
\usepackage[
  letterpaper,
  left=0.5in,
  right=0.5in,
  top=0.5in,
  bottom=1.0in
]{geometry}

% ============================
% Required packages (always)
% ============================
\usepackage[most]{tcolorbox}
\usepackage{booktabs}
\usepackage{fancyhdr}
\usepackage{graphicx}
\usepackage{lastpage}
\usepackage{listings}
\usepackage{setspace}
\usepackage{subcaption}
\usepackage{tabularx}
\usepackage{tabularray}
\usepackage{tikz}
\usepackage{xcolor}             % [table] is set via PassOptionsToPackage above
\usepackage{xparse}
\usepackage[siunitx, american, RPvoltages]{circuitikz}
\usepackage{array}
\usepackage{colortbl}
\usepackage{pdflscape}
\usepackage{pgfplots}
\usepackage{xurl}
\usetikzlibrary{arrows.meta, angles, quotes, shapes.symbols, arrows.meta}

% ============================
% minted
% ----------------------------
% In HTML mode, tex4ht has no robust support for minted: it relies on
% Pygments output that does not survive the htlatex DVI pipeline cleanly.
% We replace minted with verbatim-style fallbacks so the document still
% compiles. The PDF build keeps full minted (with -shell-escape).
% ============================
\ifSEhtml
  \usepackage{fvextra}
  % minted is NOT loaded in HTML mode. We provide a verbatim stand-in so
  % \begin{minted}[opts]{lang} ... \end{minted} still compiles. We use
  % fancyvrb's VerbatimEnvironment trick so the body is read verbatim
  % rather than tokenised by LaTeX (a +b body argument does NOT preserve
  % verbatim and breaks on '#', '%', '{', etc.).
  \newenvironment{minted}[2][]{%
    \VerbatimEnvironment
    \begin{Verbatim}[fontsize=\footnotesize]%
  }{%
    \end{Verbatim}%
  }
  % Inline replacement: \mintinline[opts]{lang}{code}
  % Note: \mintinline must accept verbatim-style content as its third arg.
  \NewDocumentCommand{\mintinline}{O{} m v}{\texttt{#3}}
  % Setting commands become no-ops in HTML mode
  \providecommand{\setminted}[2][]{}
  \providecommand{\setmintedinline}[2][]{}
\else
  \usepackage{minted}
\fi

\ifLuaTeX
  \usepackage{tagpdf}
  \tagpdfsetup{activate-all}
  \AtBeginEnvironment{tcolorbox}{\tagpdfparaOff}
  \AtEndEnvironment{tcolorbox}{\tagpdfparaOn}
\fi

% ============================
% Hyperref (load once, after most packages)
% ============================
\usepackage[
  colorlinks=true,
  linkcolor=blue,
  urlcolor=blue,
  citecolor=blue
]{hyperref}
\hypersetup{
  pdftitle={SE 423 Assignment},
  pdfauthor={Marius Juston},
}

% ============================
% unicode-math
% ----------------------------
% Only available on LuaTeX / XeTeX. NOT on pdfLaTeX. NOT on tex4ht.
% On tex4ht, attempting to load it triggers the luatexbase.sty error.
% ============================
\ifSEhtml\else
  \ifLuaTeX
    \usepackage{unicode-math}
  \else\ifXeTeX
    \usepackage{unicode-math}
  \fi\fi
\fi

\definecolor{highlight}{rgb}{0.396078431, 0, 0.721568627}

% pgfcalendar must be loaded in the preamble (here, in header.tex) rather
% than inside timline.tex. Loading \usepackage from within \input{...} is
% fragile under tex4ht (it triggers "Paragraph ended before \@reset@ptions").
\usepackage{pgfcalendar}

\input{timline}
\input{icons}

% ============================
% Assignment mode resolution
% ============================
\def\CourseCode{SE 423}
\def\CourseName{Mechatronics}

\def\AssignmentType{}
\def\AssignmentFooter{}

\makeatletter
\@namedef{Assignment@HW@type}{Homework Assignment}
\@namedef{Assignment@HW@footer}{HW}

\@namedef{Assignment@Lab@type}{Laboratory Assignment}
\@namedef{Assignment@Lab@footer}{Lab}

\@namedef{Assignment@LabVIEW@type}{LabVIEW Assignment}
\@namedef{Assignment@LabVIEW@footer}{LabVIEW}

\@ifundefined{Assignment@\AssignmentMode @type}{
  \PackageError{AssignmentMode}{Invalid AssignmentMode}{%
    Valid options: HW, Lab, LabVIEW}
}{
  \edef\AssignmentType{\@nameuse{Assignment@\AssignmentMode @type}}
  \edef\AssignmentFooter{\@nameuse{Assignment@\AssignmentMode @footer}}
}
\makeatother

% ============================
% Header / footer
% ============================
\pagestyle{fancy}
\fancyhf{}
\fancyfoot[L]{// \CourseCode, \CourseName}
\fancyfoot[R]{// \AssignmentFooter~\HWNum, Page \thepage\ of \pageref{LastPage}}
\renewcommand{\headrulewidth}{0pt}

% ============================
% Title box
% ============================
\newcommand{\MakeAssignmentTitle}{%
\par
\begin{tcolorbox}[
    colback=gray!25!white,
    width=\textwidth,
    boxrule=0.8pt,
    arc=2mm,
    left=6mm,
    right=6mm,
    top=4mm,
    bottom=4mm,
    center
]
\centering
\setstretch{1.2}

{\bfseries\Large \CourseCode\ \CourseName \\ \AssignmentType\ \#\HWNum\par}

\vspace{2mm}

{\bfseries
\AssignmentText\par
}
\end{tcolorbox}
\vspace{0.5cm}
}

% ============================
% minted styling (PDF mode only - stubs above are no-ops in HTML)
% ============================
\setmintedinline{breaklines, breakafter=_}

% Note: breaklines is intentionally OFF here. fvextra (which minted
% builds on) implements line wrapping via \everypar, and tagpdf's
% phase-III paragraph-hook counter desynchronises when a wrapped
% logical line crosses a page break — manifesting as a fatal
% "automatic begin/end text-unit para hooks differ" at \end{document}.
% Until tagpdf+fvextra cooperate cleanly under PDF/UA-2, leave wrapping
% off; long lines flow off the right margin rather than wrapping.
\setminted{
  breaklines=false,
  frame=lines,
  fontsize=\footnotesize,
}

\lstset{
    basicstyle=\ttfamily,   % Typewriter font for code
    breaklines=true,        % Allow breaking long lines
    columns=flexible        % Makes spacing better for inline
}

\newcounter{exercise}

\NewDocumentCommand{\Ex}{o}{%
  \par
  \stepcounter{exercise}
  \section*{Exercise \theexercise
    \IfValueT{#1}{: #1}
  }
  \addcontentsline{toc}{section}{Exercise \theexercise
    \IfValueT{#1}{: #1}}
}

\newcolumntype{C}[1]{>{\centering\arraybackslash}p{#1}}

\pgfplotsset{compat=1.18}

\newcommand{\MakeFooter}{%
\par
\begin{tcolorbox}[
    colback=gray!25!white,
    width=\textwidth,
    boxrule=0.8pt,
    arc=2mm,
    left=6mm,
    right=6mm,
    top=4mm,
    bottom=4mm,
    center
]
\centering
\setstretch{1.2}

{\bfseries Make sure that you \textcolor{highlight}{comment} you code and put your \textcolor{highlight}{initials} in your Gradescope submission! See \url{https://marius-juston.github.io/SE-423---Class-Material/Homeworks/commenting_code_example.pdf} for how we expect the code to be commented like!\par}
\end{tcolorbox}
}
%
\begin{document}

\MakeAssignmentTitle

\Ex[Internal commenting]

Below are examples of how we are expecting you to comment your code. \textbf{MAKE SURE YOU DO THIS!!!} Also make sure you put in the \textbf{top of the file} what the initials to be checked for are!

\begin{minted}[escapeinside=||]{c}
// Andy Jones is commenting this Code. Comments marked by my initials |\textcolor{red}{\textbf{AAJ}}|
//#############################################################################
// FILE: LABstarter_main.c
//
// TITLE: Lab Starter
//#############################################################################
#include <stdio.h>
#include <stdlib.h>
.... Omitted Code
.... You should not omit code in your submission, it has just been removed here for the sake of conciseness
....
uint32_t numTimer0calls = 0;
uint32_t numSWIcalls = 0;
extern uint32_t numRXA;
uint16_t UARTPrint = 0;
uint16_t LEDdisplaynum = 0;
float Kp = 3.4; // |\textcolor{red}{\textbf{AAJ}}| Kp and Ki gains used for the PI controller. Hand tuned for desired response.
float Ki = 29.3; // I did find that a number of gain combinations worked, but settled on these gains.
float Vref = 2.5;
float Vactual = 0;
float Ik = 0;
float Vout = 0;
void main(void)
{
    // PLL, WatchDog, enable Peripheral Clocks
    // This example function is found in the F2837xD_SysCtrl.c file.
    InitSysCtrl();
    .... Omitted Code
    .... You should not omit code in your submission, it has just been removed here for the sake of conciseness
    ....
    // Configure CPU-Timer 0, 1, and 2 to interrupt every second:
    // 200MHz CPU Freq, 1 second Period (in uSeconds)
    ConfigCpuTimer(&CpuTimer0, 200, 5000); //|\textcolor{red}{\textbf{AAJ}}| Set CPU Timer 0 to call its interrupt function every 5ms. This is the specified sample rate for the control loop calculations.
    ConfigCpuTimer(&CpuTimer1, 200, 20000);
    ConfigCpuTimer(&CpuTimer2, 200, 40000);
    .... Omitted Code
    .... You should not omit code in your submission, it has just been removed here for the sake of conciseness
    ....
    // Enable global Interrupts and higher priority real-time debug events
    EINT; // Enable Global interrupt INTM
    ERTM; // Enable Global realtime interrupt DBGM

    !!!! There is a second page !!!!!
    // IDLE loop. Just sit and loop forever
    while(1)
    {
        if (UARTPrint == 1 ) {
            // |\textcolor{red}{\textbf{AAJ}}| Here I am using GPIO67 to see how much time the serial_printf function takes to run. In lab I will connect the Oscilloscope
            // to GPIO67's pin and measure the width of the pulse created by setting GPIO67 high (3.3Volts) before the function is called and then
            // setting GPIO67 low (0Volts) after the function is completed. GPIO67 is controlled by the GPC registers. GPIO67 is the 4th bit
            // in the GPC registers. Bit 0 is for GPIO64 Bit 1 is for GPIO65, Bit 2 is for GPIO66, Bit 3 is for GPIO67, Bit 4 is for GPIO68 and so on.
            // The below bitfield C statement hides the bit number from me and allows me to just set the associated bit pertaining to GPIO67.
            // Setting a bit to 1 in the GPCSET register causes that associated output pin to be high (3.3 Volts).
            // Setting a bit to 0 in the GPCSET register does nothing.
            GpioDataRegs.GPCSET.bit.GPIO67 = 1;
            //|\textcolor{red}{\textbf{AAJ}}| Printing Desired and Actual Speed %.3f tells printf to print 3 decimal places of float variable
            serial_printf(&SerialA,"%.3f Desired Speed, %.3f Actual Speed\r\n",Vref,Vactual);
            UARTPrint = 0;
            // |\textcolor{red}{\textbf{AAJ}}| Set GPIO67 to low (0Volts) so Oscilloscope and see when serial_printf function is complete
            // Setting a bit to 1 in the GPCCLEAR register causes that associated output pin to be low (0 Volts). Setting a bit
            // to 0 in the GPCCLEAR register does nothing.
            GpioDataRegs.GPCCLEAR.bit.GPIO67 = 1
        }
    }
}
// cpu_timer0_isr - CPU Timer0 ISR
__interrupt void cpu_timer0_isr(void)
{
    Vactual = readfeedback(); // |\textcolor{red}{\textbf{AAJ}}| this is a dummy function for this example.
    // |\textcolor{red}{\textbf{AAJ}}| Here I am calculating the PI control law. I am using the "forward rule" to approximate the integral
    // Then I calculate the control effort Vout which is Kp times the error between desired speed and actual speed
    // and Ki times the integral appoximation of the error. I also check for integral windup and then saturate
    // the output between -10 and +10.
    Ik = Ikold + (Vref-Vactual)*0.005;
    Vout = Kp*(Vref-Vactual) + Ki*Ik;
    if (fabs(Vout) > 10) Ik = Ikold;
    if (Vout > 10) Vout = 10;
    if (Vout < -10) Vout = -10;
    outPWM(Vout); //|\textcolor{red}{\textbf{AAJ}}| dummy function to output control
    Ikold = Ik; //|\textcolor{red}{\textbf{AAJ}}| save past state for next 5ms
    CpuTimer0.InterruptCount++;
    numTimer0calls++;
    if ((numTimer0calls%20) == 0) {
        UARTPrint = 1; //|\textcolor{red}{\textbf{AAJ}}| Print to Tera Term every 100ms 5ms * 20 = 100ms
        // Blink LaunchPad Red LED
        GpioDataRegs.GPBTOGGLE.bit.GPIO34 = 1;
    }

    // Acknowledge this interrupt to receive more interrupts from group 1
    PieCtrlRegs.PIEACK.all = PIEACK_GROUP1;
}
.... Omitted Code
.... You should not omit code in your submission, it has just been removed here for the sake of conciseness
....
\end{minted}

\newpage
\mintinline{c}{hello}
\begin{minted}{c}
\end{minted}
%
\Ex[Function documentation]

When you specifically create a function (you do not need to do this for the provided interrupts or the main function, but you are welcome to) you should be using be using \href{https://www.doxygen.nl/manual/docblocks.html}{Doxygen JAVADOC} styling for function documentation (it is what the industry uses so use it!).

Each function should have the following structed documentation:

\begin{enumerate}
    \item Brief summary (first sentence)
    \item A more detailed description.
    \item \lstinline{@param <parameter name> <description of parameter>} for each parameter, ordered by when it appears in the function's parameters
    \item \lstinline{@return <description of what gets returned>}, when applicable (\mintinline{c}{void} functions would not have this)
    \item \lstinline{@pre <parameter expectations>}, providing expectations of what the input is supposed to satisfy when inputted into it
    \item \lstinline{@post <return expectations>}, providing expectations of how the parameters are modified during the function
    \item \lstinline{@note} or \lstinline{@warning} when relevant 
\end{enumerate}

\begin{minted}[escapeinside=||]{c}
/**
 * |\textcolor{red}{\textbf{AAJ}}|: A brief history of JavaDoc-style (C-style) comments.
 *
 * This is the typical JavaDoc-style C-style comment. It starts with two
 * asterisks.
 *
 * @param theory Even if there is only one possible unified theory. it is just a
 *               set of rules and equations.
 */
void cstyle( int theory );
\end{minted}

\begin{minted}[escapeinside=||]{c}
/**
 * |\textcolor{red}{\textbf{AAJ}}|: Computes the arithmetic mean of an array of doubles.
 *
 * @param data Pointer to the first element of the array.
 * @param length Number of elements in the array (must be > 0).
 * @return The arithmetic mean of the values in the array.
 *
 * @pre data != NULL
 * @pre length > 0
 * @post The input array is not modified.
 */
double mean(const double *data, size_t length);
\end{minted}


\begin{minted}[escapeinside=||]{c}
/**
 * |\textcolor{red}{\textbf{AAJ}}|: Opens a file and reads its contents into a buffer.
 *
 * The function allocates memory dynamically for the buffer.
 * The caller is responsible for freeing it with free().
 *
 * @param path Null-terminated string specifying the file path.
 * @param buffer Output pointer that will receive allocated memory.
 * @param size Output parameter storing the number of bytes read.
 *
 * @return 0 on success, negative error code on failure.
 *
 * @retval 0 Success.
 * @retval -1 File could not be opened.
 * @retval -2 Memory allocation failed.
 *
 * @pre path != NULL
 * @pre buffer != NULL
 * @pre size != NULL
 * @post On success, *buffer points to valid allocated memory.
 *
 * @warning The function does not append a null terminator.
 */
int read_file(const char *path, unsigned char **buffer, size_t *size);
\end{minted}

\begin{minted}[escapeinside=||]{c}
/**
 * |\textcolor{red}{\textbf{AAJ}}|: Initializes the ADC peripheral.
 *
 * Configures hardware registers and enables interrupts required
 * for analog-to-digital conversion.
 *
 * @param config Pointer to ADC configuration structure.
 *
 * @return true if initialization succeeds, false otherwise.
 *
 * @pre config != NULL
 * @post ADC hardware is ready for sampling.
 *
 * @note This function must be called before adc_read().
 */
bool adc_init(const adc_config_t *config);
\end{minted}

\Ex[Function template]

Below is a template that you can use,

\begin{minted}[escapeinside=||]{c}
/**
 * |\textcolor{red}{\textbf{AAJ}}|: Brief one-line summary.
 *
 * Detailed description if needed.
 *
 * @param name Description.
 * @return Description.
 *
 * @pre Conditions that must hold.
 * @post Guarantees after execution.
 */
\end{minted}

\end{document}